\chapter{Glomerular Filtration Rate (GFR) Test}
\section*{What Is This Test?} GFR estimates how well the kidneys filter blood and is central to chronic kidney disease assessment.\section*{Alternative Names} eGFR; estimated glomerular filtration rate; kidney filtration test.\section*{Principle} Equations use serum creatinine and sometimes cystatin C with demographic and clinical variables to estimate filtration.\section*{Why This Test Is Ordered} It screens for, stages, and monitors kidney disease and helps guide medication dosing.\section*{Primary vs Secondary Test} eGFR is a primary kidney-function measure; measured clearance or nuclear tests are confirmatory in selected cases.\section*{How This Test Is Performed} Creatinine is measured from blood and eGFR is calculated; urine albumin and other kidney tests add context.\section*{What Preparation Is Needed} Usually none. Report unusual muscle mass, amputation, pregnancy, diet, supplements, and recent acute illness.\section*{Understanding Results} Lower eGFR may indicate reduced filtration, but trends and duration matter; acute kidney injury can make estimates unreliable.\section*{Follow-up Tests} Repeat creatinine/eGFR, urine albumin-creatinine ratio, urinalysis, renal ultrasound, or nephrology referral may follow.\section*{References} \begin{enumerate}\item MedlinePlus. Glomerular Filtration Rate (GFR) Test.\item Kidney Disease: Improving Global Outcomes. CKD evaluation guidance.\end{enumerate}\clearpage\section*{Understanding Results and Follow-up}
The laboratory report should identify the specimen, method, units, reference interval or decision limit, and whether the result is positive, negative, abnormal, or inconclusive. Reference intervals vary with age, sex, pregnancy, specimen type, assay, and laboratory; a result outside a range is not automatically a diagnosis, and a result within a range does not always exclude disease. Discuss unexpected results with the ordering clinician, who can decide whether repeat or confirmatory testing, treatment, or referral is appropriate. Do not change medicines or delay urgent care based on an isolated result.
\section*{Regulatory Status and Authorization Date}This chapter describes a test category, not a specific commercial product. FDA, CDSCO, EU IVDR, and MHRA approval or registration dates therefore cannot be assigned without the assay manufacturer, product name, instrument, and intended use. For laboratory-developed tests, an individual agency approval date may not exist. See the Preface for the interpretation of regulatory status.
\clearpage
